% ===================================================================
%  CADRO N.º 14 — A rúa. — Os comerciantes.
%  Traduction 2026 d'apres texto/io, controlee sur texto/fr, texto/es
%  et texto/pt. Consigne : voir texto/gl/10-cadro-01.tex.
%
%  LE TABLEAU DES METIERS EST LE PLUS LONG DE LA COLONNE, et il donne
%  ici la meme mesure qu'en irlandais : le mot a l'age de son objet.
%  Ce qui se fait de ses mains depuis toujours garde son mot
%  galicien — « xastre » le tailleur, « peiteeiro » le coiffeur,
%  « ourive » l'orfevre, « estañador » l'etameur, « armeiro »
%  l'armurier, « folleiro » le ferblantier, « encadernador » le
%  relieur, « limpachemineas » le ramoneur — et ce qui s'est invente
%  au XIXe siecle arrive avec son nom : « fotógrafo », « litógrafo »,
%  « mecanógrafo », « estenógrafo », « óptico », « ciclista »,
%  « automóbil », « triciclo ».
%
%  ET LE GALICIEN AJOUTE UNE OBSERVATION QUE L'IRLANDAIS NE POUVAIT
%  PAS FAIRE : les mots neufs y sont tous batis sur le meme suffixe
%  grec, « -grafo », qui n'a rien de galicien et rien d'espagnol non
%  plus. Ce n'est pas une langue qui emprunte a sa voisine, ce sont
%  DEUX langues qui empruntent au meme endroit, en meme temps. La
%  ressemblance des colonnes voisines, dans ce tableau, ne prouve
%  donc aucun contact entre elles.
% ===================================================================

%%K t14-apar-1 apar
\VUsaut{4.0mm}
\VUtitre{15.0pt}{60}{CADRO N.\textsuperscript{o} 14}
\VUsaut{3.0mm}

%%K t14-tit-1 sub
\VUpk{11.4pt}{40}{A rúa. --- Os comerciantes.}
\VUsaut{3.0mm}

%%K t14-01-1 p
1. --- O meu amigo Xoán, que vive no campo, veu pasar tres días coa miña familia.
\VUgras{Ata} entón non tivera ocasión de ver unha cidade grande; por iso empregamos
todos os nosos días a pasear, e eu explícolle o que non coñece.

%%K t14-02-1 p
2. --- Primeiro fomos visitar o \VUgras{observatorio} \textsuperscript{(1)}, que lle
interesou moito. Ao volver pola \VUgras{rúa} \textsuperscript{(3)} na que están os
\VUgras{baños turcos} \textsuperscript{(2)}, atopamos un \VUgras{enterro}
\textsuperscript{(4)}. Diante do \VUgras{cortexo fúnebre} camiñaba o \VUgras{clero,}
que precedía o \VUgras{coche fúnebre,} no que estaba posto o \VUgras{cadaleito.}
Moitos amigos acompañaban a familia \VUgras{de loito.}

%%K t14-02-2 p
Despois de pasar o acompañamento, atravesamos a rúa diante da gran \VUgras{cervexaría}
\textsuperscript{(5)} na que moitos \VUgras{consumidores} bebían \VUgras{cervexa} na
terraza, resgardados pola \VUgras{marquesiña} envidrada \textsuperscript{(6)}.

%%K t14-03-1 p
3. --- Xoán quedou moi abraiado polo \VUgras{escudo} posto entre o almacén do
\VUgras{licoreiro} \textsuperscript{(7)} e o do \VUgras{vendedor de música}
\textsuperscript{(10)}. Preguntoume o seu significado; respondinlle que sinala o
\VUgras{consulado} \textsuperscript{(8)} inglés, e fíxenlle reparar no
\VUgras{cónsul} \textsuperscript{(9)} que estaba no balcón.

%%K t14-tit-2 sub
\VUpk{10.2pt}{40}{Exploración das tendas.}
\VUsaut{3.0mm}

%%K t14-04-1 p
4. --- Como é natural, paramos diante do escaparate dos \VUgras{instrumentos de
música,} \VUgras{harmonios,} \VUgras{órganos} e \VUgras{pianos}; miramos as cubertas
ilustradas das \VUgras{particións} e das \VUgras{pezas de música} para piano, e
admiramos a \VUgras{lira} de madeira dourada que se usa como letreiro.

%%K t14-05-1 p
5. --- As nubes de po feitas polos \VUgras{varredores} \textsuperscript{(11)} e polo
\VUgras{carro de varrer} \textsuperscript{(12)} obrigáronnos a pasar axiña diante dos
belos \VUgras{mobles} do \VUgras{vendedor de mobles} \textsuperscript{(13)} e diante
do escaparate de \VUgras{fogóns} e \VUgras{lámpadas} do \VUgras{folleiro}
\textsuperscript{(14)}.

%%K t14-06-1 p
6. --- Ademais, nese lugar vímonos obrigados a deixar o beirarrúa, pois estaba
ateigado de paquetes que se collían dun \VUgras{carro de transporte}
\textsuperscript{(16)} para os poñer nun \VUgras{almacén} \textsuperscript{(15)} que
acababan de alugar.

%%K t14-06-2 p
Iso impediunos ver os belos coches do \VUgras{carroceiro} \textsuperscript{(17)}, pero
non nos impediu parar para mirar o \VUgras{embalador} \textsuperscript{(19)} que facía
unha caixa para o seu veciño, o \VUgras{vendedor de bicicletas e automóbiles}
\textsuperscript{(18)}. Despois, como xa era algo tarde, volvemos á casa.

%%K t14-tit-3 sub
\VUpk{10.2pt}{40}{Paseo pola cidade.}
\VUsaut{3.0mm}

%%K t14-07-1 p
7. --- Ao día seguinte, miña nai propúxolle a Xoán que lle tirasen un retrato, e
despois encargoume que o acompañase á casa do \VUgras{fotógrafo}
\textsuperscript{(20)}. Despois do xantar fomos ao \VUgras{taller} do artista, onde
nos vimos obrigados a agardar bastante. Para nos entreter, miramos os
\VUgras{retratos} \textsuperscript{(22)} pendurados na parede.

%%K t14-07-2 p
Cando chegou a nosa quenda, o traballo non durou moito e, axiña que o meu amigo acabou
de pousar diante da \VUgras{cámara} \textsuperscript{(21)}, baixamos.

%%K t14-08-1 p
8. --- No primeiro andar vimos, pola porta do \VUgras{peiteeiro}
\textsuperscript{(23)}, que estaba aberta, o seu oficial que cortaba o pelo a un
cliente.

%%K t14-08-2 p
Ao chegar á rúa pasamos diante do \VUgras{estanco} \textsuperscript{(24)}. Xoán quedou
tan divertido co \VUgras{farol} vermello \textsuperscript{(25)} e coa \VUgras{pipa
grosa} que se usa como \VUgras{letreiro} \textsuperscript{(26)}, que chocou con dous
vellos señores que conversaban \VUgras{tomando rapé} \textsuperscript{(27)}.

%%K t14-tit-4 sub
\VUpk{10.2pt}{40}{A pastelaría.}
\VUsaut{3.0mm}

%%K t14-09-1 p
9. --- Os \VUgras{billetes de banco} e as \VUgras{moedas} de todos os países expostos
no \VUgras{escaparate} do \VUgras{cambista} \textsuperscript{(29)} atraeron a nosa
atención uns intres, pero, como era a hora de merendar, entramos na casa do
\VUgras{pasteleiro} \textsuperscript{(31)} sen parar máis tempo.

%%K t14-10-1 p
10. --- Ao ver todas as \VUgras{lambetadas} que había no almacén, \VUgras{pasteis,
pan de mel, biscoitos} e toda clase de \VUgras{caramelos,} non nos foi doado escoller.
Finalmente Xoán escolleu os \VUgras{pasteis de nata,} e eu tomei só unha pequena
\VUgras{torta de cereixas,} pois os doces \VUgras{estrágame os dentes,} e non teño
ningunha gana de ir moi a miúdo á casa do \VUgras{dentista} \textsuperscript{(30)}.

%%K t14-tit-5 sub
\VUpk{10.2pt}{40}{A escritura a máquina.}
\VUsaut{3.0mm}

%%K t14-11-1 p
11. --- Para lle amosar ao meu amigo unha \VUgras{máquina de escribir} da que oíra
falar, pero que non coñecía, entramos no \VUgras{despacho} \textsuperscript{(32)} do
meu \VUgras{curmán} \textsuperscript{(34)}. Atopámolo ditando as súas cartas ao seu
\VUgras{secretario} \textsuperscript{(35)}, que é \VUgras{estenógrafo}. Apenas
rematada unha carta, un \VUgras{mecanógrafo} \textsuperscript{(33)} copiábaa coa súa
máquina. Como non queriamos interromper as ocupacións do meu parente, quedamos na súa
casa só uns intres.

%%K t14-12-1 p
12. --- Baixo o despacho do meu curmán vive un gran \VUgras{reloxeiro-xoieiro}
\textsuperscript{(37)} que é ademais ourive. O seu espléndido almacén contén moitos
\VUgras{reloxos, reloxos de péndulo, reloxos de peto,} \VUgras{cadeíñas} e
\VUgras{xoias} preciosas, por exemplo \VUgras{colares de perlas,} \VUgras{pulseiras}
de ouro, \VUgras{pendentes} e \VUgras{aneis} adornados con \VUgras{pedras} e
\VUgras{diamantes}. Un dos escaparates está reservado especialmente á
\VUgras{ourivaría} e contén unha importante colección de \VUgras{cubertos} de prata.

%%K t14-13-1 p
13. --- Ao virar na esquina da rúa reparei en que mudaran a \VUgras{placa}
\textsuperscript{(36)} posta preto do \VUgras{número} da casa. Ía dicir en voz alta a
miña observación cando se oíron os sons ledos da \VUgras{música} militar. Botamos a
correr ao encontro do rexemento, sen pensar que \VUgras{este} había pasar diante dos
almacéns do \VUgras{sombreireiro} \textsuperscript{(39)} e do \VUgras{luveiro}
\textsuperscript{(40)}, e que estariamos igual de ben situados para ver o seu desfile
quedando diante do escaparate do \VUgras{óptico} \textsuperscript{(38)}, ben
provisto de \VUgras{lentes de pinza, lentes de aro} e \VUgras{binóculos}.

%%K t14-tit-6 sub
\VUpk{10.2pt}{40}{Desfile.}
\VUsaut{3.0mm}

%%K t14-14-1 p
14. --- Diante do \VUgras{rexemento} \textsuperscript{(41)} camiñaba o
\VUgras{tambor maior} \textsuperscript{(42)}, seguido polos \VUgras{tambores}
\textsuperscript{(43)}, os \VUgras{corneteiros} \textsuperscript{(44)} e toda a
\VUgras{banda de música}. O \VUgras{xeneral} \textsuperscript{(45)}, que estaba na
praza co seu axudante, parara preto do \VUgras{chorro de auga}
\textsuperscript{(49)} da \VUgras{fonte} \textsuperscript{(48)} para asistir ao
\VUgras{desfile.}

%%K t14-14-2 p
\VUgras{Os oficiais do estado maior} \textsuperscript{(46)}, o \VUgras{coronel} e o
\VUgras{tenente coronel} saudárono coa espada. \VUgras{Os comandantes} estaban diante
dos seus \VUgras{batallóns} \textsuperscript{(47)} e cada \VUgras{capitán} mandaba a
súa \VUgras{compañía,} mentres os \VUgras{tenentes,} \VUgras{alféreces} e
\VUgras{suboficiais} ocupaban o seu sitio de costume ao lado dos seus homes.

%%K t14-15-1 p
15. --- Moitos transeúntes agrupáranse no \VUgras{cruzamento}
\textsuperscript{(50)}; os comerciantes, o \VUgras{paraugueiro} \textsuperscript{(51)},
o \VUgras{estañador} \textsuperscript{(53)} e o \VUgras{armeiro}
\textsuperscript{(54)} saíron para ver os \VUgras{soldados.} Todos os vehículos,
\VUgras{coches de aluguer} \textsuperscript{(52)}, \VUgras{carros de mudanzas}
\textsuperscript{(57)} e mesmo a \VUgras{cuba de rega} \textsuperscript{(56)},
puxéronse ao carón do beirarrúa.

%%K t14-tit-7 sub
\VUpk{10.2pt}{40}{Escenas na rúa.}
\VUsaut{3.0mm}

%%K t14-15-2 p
Un \VUgras{viaxante de comercio} \textsuperscript{(55)} que levaba na man as súas
\VUgras{caixas de mostras} atravesou axiña a rúa, e as persoas sentadas na
\VUgras{imperial} \textsuperscript{(59)} do \VUgras{ómnibus} \textsuperscript{(58)}
volvéronse para gozar do espectáculo. Un pobre \VUgras{cantor ambulante}
\textsuperscript{(60)} co seu \VUgras{organiño} \textsuperscript{(61)} tivo que
interromper a súa \VUgras{canción,} pois o ruído dos tambores cubría a súa voz.

%%K t14-16-1 p
16. --- Cando o rexemento chegou diante do \VUgras{comercio de tecidos}
\textsuperscript{(64)}, as \VUgras{costureiras} \textsuperscript{(63)} e os
\VUgras{xastres} \textsuperscript{(62)} deixaron as súas \VUgras{máquinas de coser} e
o seu traballo para mirar pola fiestra. \VUgras{O comerciante}
\textsuperscript{(65)}, que acababa de pór \VUgras{traxes} novos
\textsuperscript{(66)} no seu \VUgras{escaparate,} tivo que facer meter dentro os
\VUgras{panos} \textsuperscript{(67)} e as \VUgras{teas} postos no beirarrúa, preto do
\VUgras{sumidoiro} \textsuperscript{(68)}, por medo de que a multitude os derrubase;
mesmo un vello \VUgras{vendedor ambulante} \textsuperscript{(69)}, ao que unha
porteira lle regateaba unhas medias, viuse obrigado a entrar nun corredor para rematar
a súa venda.

%%K t14-16-2 p
Acompañamos o rexemento ata o cuartel \VUgras{e,} moi contentos coa nosa xornada,
volvemos á hora da cea.

%%K t14-tit-8 sub
\VUpk{10.2pt}{40}{Comercios e oficios diversos.}
\VUsaut{3.0mm}

%%K t14-17-1 p
17. --- Ao día seguinte, meu pai propúxonos visitar a imprenta do xornal de alí.
Aceptamos con entusiasmo.

%%K t14-17-2 p
O \VUgras{impresor} é tamén \VUgras{libreiro} \textsuperscript{(70)} \VUgras{(=
vendedor de libros),} \VUgras{editor,} \VUgras{papeleiro} e \VUgras{encadernador}.
Instalou no primeiro andar a \VUgras{redacción} \textsuperscript{(71)} do xornal, e
tamén o \VUgras{taller de gravado,} no que traballan os \VUgras{litógrafos}
\textsuperscript{(72)} e os \VUgras{gravadores en cobre} \textsuperscript{(73)}, e
finalmente o \VUgras{taller de composición,} no que están os \VUgras{obreiros
impresores} \textsuperscript{(74)}. A \VUgras{imprenta} propiamente dita
\textsuperscript{(75)}, as \VUgras{prensas de imprimir} e as distintas máquinas están
no andar baixo.

%%K t14-tit-9 sub
\VUpk{10.2pt}{40}{Xoán fai moitísimas preguntas.}
\VUsaut{3.0mm}

%%K t14-18-1 p
18. --- Ao saír da imprenta, Xoán, moi abraiado, parou diante dunha pequena caixa
envidrada posta nunha columna, contra a \VUgras{mercería} \textsuperscript{(77)}, e
preguntou para que servía. Meu pai explicoulle que é un \VUgras{avisador de
incendios} \textsuperscript{(76)}. Ademais, todo na cidade era, para o meu amigo,
motivo de abraio, desde a simple \VUgras{fonte de pedra} \textsuperscript{(78)} e o
lixeiro \VUgras{triciclo de reparto} \textsuperscript{(80)} guiado por un
\VUgras{botones} \textsuperscript{(79)} de librea, ata o \VUgras{coche de correos}
\textsuperscript{(81)}.

%%K t14-19-1 p
19. --- Mentres meu pai nos deixaba uns intres para falar co \VUgras{recadador de
impostos} \textsuperscript{(83)}, que parara para conversar cun \VUgras{avogado}
\textsuperscript{(82)} e cun \VUgras{notario} \textsuperscript{(84)}, nós
entretivémonos seguindo cos ollos o movemento da rúa. As persoas máis diversas
pasaban diante de nós: un \VUgras{oficial de pasteleiro} \textsuperscript{(85)}, que
levaba nun cesto un \VUgras{pastelón} espléndido, viña detrás dun \VUgras{vendedor de
roupa} \textsuperscript{(86)}; un \VUgras{acendedor de farois}
\textsuperscript{(87)} conversaba cun \VUgras{gasista} \textsuperscript{(88)}; unha
\VUgras{relixiosa} \textsuperscript{(89)} que levaba da man unha orfa volvía ao seu
convento.

%%K t14-tit-10 sub
\VUpk{10.2pt}{40}{No noso regreso á casa.}
\VUsaut{3.0mm}

%%K t14-20-1 p
20. --- No intre en que meu pai nos alcanzou, Xoán botou a rir con forza ao ver a cara
sucia e a roupa estraña de dous \VUgras{limpachemineas}
\textsuperscript{(91)} totalmente cheos de feluxe.

%%K t14-21-1 p
21. --- Eu quería mercarlle unha planta á miña nai á \VUgras{florista}
\textsuperscript{(92)}, pero meu pai fíxome reparar en que sería moi pesada de levar e
que sería mellor mercar \VUgras{laranxas}. Achegámonos á \VUgras{vendedora}
\textsuperscript{(94)} e, cando a \VUgras{institutriz} \textsuperscript{(93)}, que
estaba a escoller para a súa pupila, rematou a súa compra, fixemos que nos servisen.

%%K t14-22-1 p
22. --- Ao volver á casa atopamos varios coñecidos: unha vella amiga da miña familia,
que apoiaba o brazo na súa \VUgras{dama de compaña} \textsuperscript{(95)}, o
\VUgras{enxeñeiro} \textsuperscript{(96)} de pontes e camiños, e unha das miñas
curmás coa súa \VUgras{neneira} \textsuperscript{(98)}.

%%K t14-22-2 p
Despois atopamos a \VUgras{modista} \textsuperscript{(97)} da miña nai e dous
\VUgras{monxes} \textsuperscript{(99)} forasteiros na cidade, que se achegaron a nós
para lle preguntar a meu pai polo camiño da estación.
